\subsubsection{6.2 预期研究结果}

{\kaishu
	本项目力争在海杂波场景电磁频谱数据智能处理机理与方法研究方面取得一系列原创性成果，以奠定项目团队在该方向上的国际领先学术地位。预期的研究成果将以技术报告、论文、专利、软著等形式给出，概括如下：
	
	\textbf{（1）理论和技术方面：}
	深入研究海杂波场景电磁频谱数据智能处理关键科学问题，提出海杂波干扰下微弱电磁频谱信号传播规律与微弱信号提取规律、建立一套针对频谱数据多维表征的标准、有限算力端侧频谱数据的轻量存储和可信管理方法、以及基于频谱知识图谱的轻量化目标识别与可信协同决策方法。
	
	\textbf{（2）论文方面：}
	在电磁频谱智能处理、人工智能、数据库和区块链等相关领域具有国际影响力的国内科技期刊（如《中国科学：信息科学》、《计算机学报》、《软件学报》等）、业界公认的国际顶级或重要科技期刊（如 TSP、TGRS、TKDE、TIFS 等）以及国内外顶级学术会议（如 NeurIPS、ICLR、SIGMOD、VLDB、ICDE、ICASSP 等）上发表/录用25篇“三类高质量论文”，EI会议论文15篇。
	
	%\textbf{（3）应用验证方面：}
	%集成所取得的研究成果，开展云物融合架构下人机交互的具身智能感知机理与协同演化研究。在智慧管廊巡检场景下进行原型系统应用验证，提供气体超标、热源异常、水位上涨等典型应急用例的全链路闭环处置，验证系统在通信断联与资源波动约束下的任务连续性与执行效率，推动关键技术向标准化巡检业务的转化应用。
	
	\textbf{（3）专利软著方面：}
	申请3项相关的国家/国际发明专利和3项软件著作权，并注重专利落地转化，突出产业化应用的实际效果。
	
	\textbf{（4）人才培养方面：}
	依托现有的项目团队，培养博士研究生10人和硕士研究生10人（争取培养行业学会优秀博士学位论文奖获得者1人、国际知名学术会议最佳或优秀论文获得者1人），形成具有较强影响力的研究团队，进一步提升国内外的学术研究地位，在海杂波场景电磁频谱数据智能处理等前沿领域取得高水平的创新成果。
	
	\textbf{（5）学术交流方面：}
	定期组织与本项目科研任务相关的学术研讨会，派遣项目组成员积极参加国内外相关的学术会议，以进行学术交流和成果展示，并邀请国内外同行专家开展学术合作交流。
}
\begin{table}[htbp]
	\kaishu
	\centering
	\caption{预期研究成果}
	\label{tab:Achievement-plan}
	\scriptsize
	\setlength{\tabcolsep}{3pt}
	\renewcommand{\arraystretch}{1.35}

	\begin{tabular}{
		|>{\centering\arraybackslash}p{0.3\textwidth}
		|>{\raggedright\arraybackslash}p{0.6\textwidth}|
	}
		\hline
		\textbf{成果类别} &
		\centering\arraybackslash\textbf{具体预期成果与指标}
%		\centering\arraybackslash\textbf{理论和技术方面} &
%		\centering\arraybackslash\textbf{论文方面} &
%		\centering\arraybackslash\textbf{专利软著方面} &
%		\centering\arraybackslash\textbf{人才培养方面} &
%		\centering\arraybackslash\textbf{学术交流方面}
		\\
		\hline
		\kaishu
		\textbf{理论和技术方面} &
		1. 提出海杂波干扰下微弱电磁频谱信号传播规律与微弱信号提取规律； \newline
		2. 建立一套针对频谱数据多维表征的标准； \newline
		3. 建立有限算力端侧频谱数据的轻量存储和可信管理方法； \newline
		4. 提出基于频谱知识图谱的轻量化目标识别与可信协同决策方法。
		\\
		\hline

		\textbf{论文方面} &

		1. 发表/录用25篇“三类高质量论文”（涵盖国内科技期刊、国际顶级或重要科技期刊及国内外顶级学术会议）； \newline
		2. 发表/录用15篇EI会议论文。

		\\
		\hline

		\textbf{专利软著方面} &

		1. 申请3项相关的国家/国际发明专利； \newline
		2. 申请3项软件著作权； \newline
		3. 注重专利落地转化，突出产业化应用的实际效果。
		\\
		\hline

		\textbf{人才培养方面} &


		1. 培养博士研究生10人和硕士研究生10人； \newline
		2. 争取培养行业学会优秀博士学位论文奖获得者1人、国际知名学术会议最佳或优秀论文获得者1人； \newline
		3. 形成具有较强影响力的研究团队。 \\
		\hline

		\textbf{学术交流方面} &

		1. 定期组织与本项目科研任务相关的学术研讨会； \newline
		2. 派遣项目组成员积极参加国内外相关的学术会议进行学术交流和成果展示； \newline
		3. 邀请国内外同行专家开展学术合作交流。 \\
		\hline

	\end{tabular}
\end{table}